\section*{Chapitre \ref{ch:b2:fragment-orbitals} --- Orbitales de fragments et molécules polyatomiques}

\begin{solution}{exo:b2:fragment-orbitals:1}
$h_1 + h_2$ et $h_1 - h_2$. Le plan $xz$, qui contient la bissectrice de l'angle et est
perpendiculaire à la molécule, échange les deux hydrogènes~: $h_1 - h_2$ change
de signe, $h_1 + h_2$ non.
\end{solution}

\begin{solution}{exo:b2:fragment-orbitals:2}
Six orbitales de valence (2s et trois 2p de l'oxygène, deux 1s des hydrogènes), huit électrons
de valence, quatre orbitales remplies.
\end{solution}

\begin{solution}{exo:b2:fragment-orbitals:3}
Les niveaux de valence remplis sont une orbitale $a_1$ et un ensemble de trois orbitales $t_2$
dégénérées~: deux énergies, deux bandes. La bande $t_2$ arrache des électrons à trois
orbitales contenant six électrons, contre une orbitale et deux électrons pour $a_1$.
\end{solution}

\begin{solution}{exo:b2:fragment-orbitals:4}
La liaison $\pi$ exige que les deux orbitales p des carbones soient parallèles, ce qui place
les deux plans \ce{CH2} dans un même plan.
\end{solution}

\begin{solution}{exo:b2:fragment-orbitals:5}
L'ammoniac perd un électron de son doublet non liant, une orbitale essentiellement non liante localisée sur
l'azote~; le méthane, d'une \omterm{def:b2:diatomic-mos:bonding}{orbitale liante} $t_2$. Un électron non liant se situe plus haut,
plus près du niveau atomique, qu'un électron liant~: il est plus facile à arracher.
\end{solution}

\begin{solution}{exo:b2:fragment-orbitals:6}
L'eau linéaire possède deux orbitales p non liantes dégénérées et remplies~; la courbure abaisse l'une
d'elles et élève légèrement un niveau liant~: gain net avec huit électrons. \ce{BeH2}
n'en a que quatre, qui remplissent les deux \omterm{def:b2:diatomic-mos:bonding}{orbitales liantes}~; l'orbitale que la courbure
abaisserait est vide, tandis que les niveaux liants perdent du recouvrement~: la forme linéaire est la meilleure.
\end{solution}

\begin{solution}{exo:b2:fragment-orbitals:7}
Chaque liaison C--H ou C--C portée par un carbone voisin du centre cationique peut s'aligner avec
l'orbitale p vacante et apporte une stabilisation d'environ $2\beta^2/\Delta$. \ce{CH3+} n'a aucun voisin
de ce type~; les cations éthyle, isopropyle et tert-butyle ont un, deux et trois groupes
méthyle voisins du centre, chacun offrant une liaison alignée~: la stabilité croît dans cet
ordre.
\end{solution}

\begin{solution}{exo:b2:fragment-orbitals:8}
Pour deux niveaux égaux, la séparation vaut $2|\beta|$ et les deux électrons $\pi$ gagnent $2|\beta|$,
proportionnel à $\cos\theta$. Ce gain tombe à la moitié pour $\theta = 60^\circ$ et s'annule à
$90^\circ$.
\end{solution}

\begin{solution}{exo:b2:fragment-orbitals:9}
Le borane possède six électrons de valence~: ils remplissent les trois \omterm{def:b2:diatomic-mos:bonding}{orbitales liantes}, et l'orbitale p
du bore perpendiculaire au plan reste vide~; la pyramidalisation n'apporte rien.
L'ammoniac en a deux de plus~: ils occupent l'orbitale qui est non liante dans
la forme plane et, comme dans l'eau, la pyramidalisation lui permet de se mélanger avec l'orbitale 2s et
la combinaison des hydrogènes, et de descendre.
\end{solution}

\begin{solution}{exo:b2:fragment-orbitals:10}
Les équations séculaires pour $c_1 = c_2 = c_3$ donnent $\alpha + 2\beta$~; les deux combinaisons
qui lui sont orthogonales donnent $\alpha - \beta$ (dégénérées). Deux électrons dans $\alpha + 2\beta$~: énergie
$2\alpha + 4\beta$, inférieure à celle de $\ce{H2} + \ce{H+}$ ($2\alpha + 2\beta$) puisque $\beta < 0$.
\end{solution}

\begin{solution}{exo:b2:fragment-orbitals:11}
$\pi_1$~: les trois coefficients de même signe, aucun nœud (liante)~; $\pi_2$~: signes opposés
sur les carbones terminaux et un nœud au centre (non liante)~; $\pi_3$~: signes
alternés, deux nœuds (antiliante). L'anion possède quatre électrons $\pi$~: $\pi_2$ est son
plus haut niveau occupé, avec une densité sur les seuls carbones terminaux.
\end{solution}

\begin{solution}{exo:b2:fragment-orbitals:12}
Dans chacun des deux plans contenant l'axe, trois orbitales p parallèles (O, C, O)
donnent une \omterm{def:b2:diatomic-mos:bonding}{orbitale liante}, une non liante (sur les seuls oxygènes) et une antiliante~: six
orbitales $\pi$ en tout. Les deux liantes et les deux non liantes sont remplies~: huit
électrons $\pi$ (deux liaisons $\pi$ et deux doublets non liants des oxygènes dans la représentation de Lewis).
\end{solution}

\begin{solution}{pb:b2:fragment-orbitals:1}
\textbf{1.} Huit~: six de l'oxygène, un de chaque hydrogène.
\textbf{2.} $h_1 + h_2$ et $h_1 - h_2$.
\textbf{3.} 2s et $2\mathrm p_z$.
\textbf{4.} $2\mathrm p_y$.
\textbf{5.} $2\mathrm p_x$, perpendiculaire au plan moléculaire.
\textbf{6.} Six~; quatre remplies.
\textbf{7.} $h_1 - h_2$, qui change de signe, comme $2\mathrm p_z$, par réflexion dans le plan
perpendiculaire à l'axe passant par l'oxygène.
\textbf{8.} $h_1 + h_2$.
\textbf{9.} $2\mathrm p_x$ et $2\mathrm p_y$, perpendiculaires à l'axe~: une paire non liante
dégénérée.
\textbf{10.} (2s, $h_1 + h_2$) liante$^2$, ($2\mathrm p_z$, $h_1 - h_2$) liante$^2$, puis
$2\mathrm p_x^2\,2\mathrm p_y^2$.
\textbf{11.} Non~: la paire dégénérée est pleine, tous les électrons sont appariés.
\textbf{12.} À l'énergie d'une orbitale 2p de l'oxygène~: il est non liant.
\textbf{13.} $2\mathrm p_y$ (l'axe de la molécule linéaire étant $z$, courbure dans $yz$).
\textbf{14.} Il se mélange avec 2s et $h_1 + h_2$ et descend~: c'est le niveau $3a_1$.
\textbf{15.} Le niveau liant ($2\mathrm p_z$, $h_1 - h_2$), qui perd un peu de recouvrement quand les
hydrogènes quittent l'axe~: il devient $1b_2$.
\textbf{16.} Le niveau rempli qui descend gagne plus que l'autre ne perd~: avec huit
électrons, la forme coudée est la plus stable.
\textbf{17.} $104.48^\circ$, entre $90^\circ$ (liaisons formées d'orbitales p seules) et $109.5^\circ$~:
l'orbitale 2s prend part à la liaison.
\textbf{18.} $2\mathrm p_x$, qui reste l'\omterm{def:b2:diatomic-mos:bonding}{orbitale non liante} $1b_1$.
\textbf{19.} Quatre.
\textbf{20.} De $1b_1$, l'orbitale $2\mathrm p_x$ non liante de l'oxygène.
\textbf{21.} Arracher un électron non liant modifie à peine les liaisons~: l'ion a la
géométrie de la molécule et peu d'énergie vibrationnelle est excitée. Arracher un
électron liant modifie les longueurs de liaison et les angles, et la bande s'étale sur
de nombreux niveaux vibrationnels.
\textbf{22.} \qty{12.62}{eV} contre \qty{13.62}{eV}~: l'oxygène de l'eau porte une charge partielle
négative prise aux hydrogènes, qui élève ses niveaux 2p.
\textbf{23.} Ce ne sont pas deux niveaux égaux~: leurs deux combinaisons sont les orbitales $3a_1$ et
$1b_1$, d'énergies différentes. Les deux représentations décrivent la même densité électronique
totale~; seules les orbitales délocalisées donnent les énergies d'ionisation.
\textbf{24.} L'eau présente \textbf{quatre} bandes d'ionisation de valence, la plus basse vers
$\boldsymbol{\approx \qty{12.6}{eV}}$.
\end{solution}
